\section{Entangled or Product; Unitary or Not}\label{sec:entangled}

\paragraph{The test for two qubits.}
Write the state as $a\ket{00} + b\ket{01} + c\ket{10} + d\ket{11}$. Suppose it is a product
$\left(\alpha_0\ket{0} + \alpha_1\ket{1}\right)\left(\beta_0\ket{0} + \beta_1\ket{1}\right)$.
Multiplying out gives $a = \alpha_0\beta_0$, $b = \alpha_0\beta_1$, $c = \alpha_1\beta_0$,
$d = \alpha_1\beta_1$, and therefore
\[ad = \alpha_0\alpha_1\beta_0\beta_1 = bc.\]
So if $ad \neq bc$ the state is entangled. This is the argument for the Bell state (Sep 15) written
once for all states. To show that a state is a product, write down the two factors.

\begin{exercisebox}
\begin{pproblem}[Five states]\label{pp:five}
    Which of these states are entangled? Prove each answer.
    \begin{enumerate}[label=(\alph*)]
        \item $\frac{1}{2}\left(\ket{00} + \ket{01} + \ket{10} + \ket{11}\right)$
        \item $\frac{1}{2}\left(\ket{00} - \ket{01} + \ket{10} - \ket{11}\right)$
        \item $\frac{1}{2}\left(\ket{00} + \ket{01} + \ket{10} - \ket{11}\right)$
        \item $\frac{1}{\sqrt{2}}\left(\ket{01} - \ket{10}\right)$
        \item $\sqrt{\sfrac{1}{3}}\, \ket{00} + \sqrt{\sfrac{2}{3}}\, \ket{01}$
    \end{enumerate}
\end{pproblem}
\tcblower
\begin{wsolution}
    Compare $ad$ with $bc$:
    \begin{center}
        \begin{tabular}{cccl}
            \toprule
            State & $ad$              & $bc$              & Answer                                                                      \\ \midrule
            (a)   & $\sfrac{1}{4}$    & $\sfrac{1}{4}$    & Product: $\ket{+}\ket{+}$                                                   \\
            (b)   & $-\sfrac{1}{4}$   & $-\sfrac{1}{4}$   & Product: $\ket{+}\ket{-}$                                                   \\
            (c)   & $-\sfrac{1}{4}$   & $\sfrac{1}{4}$    & Entangled                                                                   \\
            (d)   & $0$               & $-\sfrac{1}{2}$   & Entangled                                                                   \\
            (e)   & $0$               & $0$               & Product: $\ket{0}\left(\sqrt{\sfrac{1}{3}}\,\ket{0} + \sqrt{\sfrac{2}{3}}\,\ket{1}\right)$ \\ \bottomrule
        \end{tabular}
    \end{center}

    \emph{Full write-up for (c).} Suppose the state equals
    $\left(\alpha_0\ket{0} + \alpha_1\ket{1}\right)\left(\beta_0\ket{0} + \beta_1\ket{1}\right)$.
    Then $\alpha_0\beta_0 = \frac{1}{2}$, $\alpha_0\beta_1 = \frac{1}{2}$,
    $\alpha_1\beta_0 = \frac{1}{2}$ and $\alpha_1\beta_1 = -\frac{1}{2}$. The product of the first
    and the last is $\alpha_0\alpha_1\beta_0\beta_1 = -\frac{1}{4}$. The product of the middle two
    is $\alpha_0\alpha_1\beta_0\beta_1 = +\frac{1}{4}$. This is a contradiction, so the state is
    entangled.
\end{wsolution}
\end{exercisebox}

\begin{exercisebox}
\begin{pproblem}[Is the matrix unitary?]\label{pp:unitary}
    Decide for each matrix whether it is unitary.
    \[U = \frac{1}{\sqrt{2}} \begin{pmatrix} 1 & i \\ i & 1 \end{pmatrix},
        \qquad
        V = \begin{pmatrix} 1 & 1 \\ 0 & 1 \end{pmatrix},
        \qquad
        W = \frac{1}{\sqrt{2}} \begin{pmatrix} 1 & 1 \\ 1 & 1 \end{pmatrix}.\]
\end{pproblem}
\tcblower
\begin{wsolution}
    Compute $U^\dagger U$, where $\dagger$ means: transpose, and conjugate every entry.
    \[U^\dagger U = \frac{1}{2} \begin{pmatrix} 1 & -i \\ -i & 1 \end{pmatrix} \begin{pmatrix} 1 & i \\ i & 1 \end{pmatrix}
        = \frac{1}{2} \begin{pmatrix} 1 - i^2 & i - i \\ -i + i & -i^2 + 1 \end{pmatrix}
        = \begin{pmatrix} 1 & 0 \\ 0 & 1 \end{pmatrix}.\]
    So $U$ is unitary. Do not forget the conjugate: without it the diagonal would be $1 + i^2 = 0$.

    $V$ is not unitary: $V\ket{1} = \ket{0} + \ket{1}$ has norm $\sqrt{2}$, and a unitary matrix
    preserves norms (Sep 15). $W$ is not unitary: $W\ket{0} = W\ket{1} = \ket{+}$, so two orthogonal
    states are sent to the same state; also $W\left(\ket{0} - \ket{1}\right) = 0$.

    \emph{Quick test.} A matrix is unitary exactly when its columns are orthonormal, with the
    complex inner product.
\end{wsolution}
\end{exercisebox}

\begin{exercisebox}
\begin{pproblem}[A small circuit]\label{pp:circuit}
    Apply $H$ to the first qubit and then CNOT (the first qubit is the control) to $\ket{10}$. Is
    the result entangled?
\end{pproblem}
\tcblower
\begin{wsolution}
    \leavevmode
    \begin{enumerate}
        \item $H$ on the first qubit: $\ket{10}$ becomes
              $\frac{1}{\sqrt{2}}\left(\ket{0} - \ket{1}\right) \otimes \ket{0}
                  = \frac{1}{\sqrt{2}}\left(\ket{00} - \ket{10}\right)$.
        \item CNOT flips the second qubit when the first is $1$:
              $\frac{1}{\sqrt{2}}\left(\ket{00} - \ket{11}\right)$.
        \item Test: $ad = -\frac{1}{2}$ and $bc = 0$, so the state is entangled.
    \end{enumerate}
    The same circuit sends $\ket{00}$ to $\frac{1}{\sqrt{2}}\left(\ket{00} + \ket{11}\right)$,
    $\ket{01}$ to $\frac{1}{\sqrt{2}}\left(\ket{01} + \ket{10}\right)$ and $\ket{11}$ to
    $\frac{1}{\sqrt{2}}\left(\ket{01} - \ket{10}\right)$.
\end{wsolution}
\end{exercisebox}

\section{Random Walk on a Small Graph}\label{sec:walk}

The recipe (Sep 8, Sep 10), for a graph with $m$ edges:
\[P_{uv} = \begin{cases} \dfrac{1}{\deg(u)} & \text{if } u \sim v, \\[1ex] 0 & \text{otherwise,} \end{cases}
    \qquad
    p_{t+1} = p_t P,
    \qquad
    \pi_v = \frac{\deg(v)}{2m}.\]
The walk converges to $\pi$ from every start exactly when the graph is connected and not bipartite.
A graph is bipartite exactly when it has no odd cycle.

\begin{exercisebox}
\begin{pproblem}[A triangle with a tail]\label{pp:triangle}
    Vertices $1, 2, 3$ form a triangle and vertex $4$ is joined to vertex $3$.
    \begin{enumerate}[label=(\alph*)]
        \item Write $P$.
        \item Find $\pi$ and verify it.
        \item Does the walk converge?
        \item Starting at vertex $4$, compute $p_1, p_2, p_3$.
    \end{enumerate}
\end{pproblem}
\tcblower
\begin{wsolution}
    (a) The degrees are $2, 2, 3, 1$, so
    \[P = \begin{pmatrix}
            0            & \tfrac{1}{2} & \tfrac{1}{2} & 0            \\
            \tfrac{1}{2} & 0            & \tfrac{1}{2} & 0            \\
            \tfrac{1}{3} & \tfrac{1}{3} & 0            & \tfrac{1}{3} \\
            0            & 0            & 1            & 0
        \end{pmatrix}.\]
    Every row sums to $1$.

    (b) There are $m = 4$ edges, so
    $\pi = \left(\frac{2}{8}, \frac{2}{8}, \frac{3}{8}, \frac{1}{8}\right)
        = \left(\frac{1}{4}, \frac{1}{4}, \frac{3}{8}, \frac{1}{8}\right)$.
    Verify $\pi P = \pi$ entry by entry:
    \begin{itemize}
        \item Vertex $1$: $\pi_2 \cdot \frac{1}{2} + \pi_3 \cdot \frac{1}{3} = \frac{1}{8} + \frac{1}{8} = \frac{1}{4}$.
        \item Vertex $2$: $\pi_1 \cdot \frac{1}{2} + \pi_3 \cdot \frac{1}{3} = \frac{1}{8} + \frac{1}{8} = \frac{1}{4}$.
        \item Vertex $3$: $\pi_1 \cdot \frac{1}{2} + \pi_2 \cdot \frac{1}{2} + \pi_4 \cdot 1
                  = \frac{1}{8} + \frac{1}{8} + \frac{1}{8} = \frac{3}{8}$.
        \item Vertex $4$: $\pi_3 \cdot \frac{1}{3} = \frac{1}{8}$.
    \end{itemize}

    (c) Yes. The graph is connected, and it contains a triangle (an odd cycle), so it is not
    bipartite.

    (d) $p_0 = (0, 0, 0, 1)$.
    \begin{itemize}
        \item $p_1 = (0, 0, 1, 0)$: from $4$ the walk must go to $3$.
        \item $p_2 = \left(\frac{1}{3}, \frac{1}{3}, 0, \frac{1}{3}\right)$: from $3$ each
              neighbour has probability $\frac{1}{3}$.
        \item $p_3 = \left(\frac{1}{3} \cdot \frac{1}{2},\ \frac{1}{3} \cdot \frac{1}{2},\
                  \frac{1}{3} \cdot \frac{1}{2} + \frac{1}{3} \cdot \frac{1}{2} + \frac{1}{3} \cdot 1,\ 0\right)
                  = \left(\frac{1}{6}, \frac{1}{6}, \frac{2}{3}, 0\right)$.
    \end{itemize}
    Each $p_t$ sums to $1$, as it must (a row stochastic matrix preserves probability vectors).
\end{wsolution}
\end{exercisebox}

\begin{exercisebox}
\begin{pproblem}[A path on three vertices]\label{pp:path}
    Consider the path $1 - 2 - 3$. Find $\pi$. Does the walk from vertex $1$ converge?
\end{pproblem}
\tcblower
\begin{wsolution}
    \leavevmode
    \begin{enumerate}
        \item The degrees are $1, 2, 1$ and $m = 2$, so
              $\pi = \left(\frac{1}{4}, \frac{1}{2}, \frac{1}{4}\right)$. Check:
              $\pi P = \left(\frac{1}{2} \cdot \frac{1}{2},\ \frac{1}{4} + \frac{1}{4},\ \frac{1}{2} \cdot \frac{1}{2}\right) = \pi$.
        \item The graph is bipartite, with sides $\set{1, 3}$ and $\set{2}$.
        \item From $p_0 = (1, 0, 0)$: $p_1 = (0, 1, 0)$,
              $p_2 = \left(\frac{1}{2}, 0, \frac{1}{2}\right)$, $p_3 = (0, 1, 0)$, and so on.
        \item The distribution alternates forever between $(0, 1, 0)$ and
              $\left(\frac{1}{2}, 0, \frac{1}{2}\right)$. It does not converge, although $\pi$
              exists and satisfies $\pi P = \pi$.
    \end{enumerate}
    In terms of eigenvalues: $P$ has the eigenvalues $1$, $0$ and $-1$. The eigenvalue $-1$ is what a
    bipartite graph produces, and its contribution $(-1)^t$ never dies out.
\end{wsolution}
\end{exercisebox}

\begin{exercisebox}
\begin{pproblem}[Short questions]\label{pp:walkshort}
    Find the stationary distribution and say whether the walk converges for
    \begin{enumerate}[label=(\alph*)]
        \item a $4$-cycle and a $5$-cycle,
        \item any $d$-regular graph on $n$ vertices,
        \item a disconnected graph.
    \end{enumerate}
\end{pproblem}
\tcblower
\begin{wsolution}
    (a) In a $4$-cycle every degree is $2$, so $\pi$ is uniform,
    $\left(\frac{1}{4}, \frac{1}{4}, \frac{1}{4}, \frac{1}{4}\right)$. An even cycle is bipartite, so
    the walk from a single vertex does not converge. A $5$-cycle is an odd cycle, so the walk does
    converge, to the uniform distribution.

    (b) $2m = nd$, so $\pi_v = \sfrac{d}{(nd)} = \sfrac{1}{n}$: the stationary distribution is
    uniform.

    (c) Each component has its own stationary distribution, so $\pi$ is not unique and the limit
    depends on the start.
\end{wsolution}
\end{exercisebox}
